% GENERATED FILE. Edit canonical lessons, then run npm run generate:books.
% canonical-source-hash: fnv1a64:da4bab1ad344b522
% canonical-lessons: GU-C416-kendra, GU-C416-vistar, GU-C416-paru, GU-C416-bhagol, GU-C416-bhamvu, GU-R416-review-409-416

\chapter{Areas and Places}
\label{ch:gu-areas-and-places}
\hlchaptermodality{416}

\begin{chapteropening}
\textbf{By the end of this chapter:} \emph{I can talk about the parts of a town and its edges.}

Complete the last lesson of chapter 416: I can talk about the parts of a town and its edges.
\end{chapteropening}

\section[\texorpdfstring{\gu{કેન્દ્ર} (kendra)}{કેન્દ્ર (kendra)}]{\gu{કેન્દ્ર} (kendra) — a centre}

\label{lesson:GU-C416-kendra}



\begin{quote}
Before the new one: say the Gujarati for a clock tower, then the Gujarati for to hunt for.
\end{quote}

\subsection*{You'll want to know: \gu{કેન્દ્ર}}
\textbf{\gu{કેન્દ્ર}} — \emph{kendra} — \textquotedblleft{}a centre\textquotedblright{}.

\begin{quote}
\gu{શહેરના} \gu{કેન્દ્રમાં} \gu{મોટું} \gu{બજાર} \gu{છે}.
\end{quote}

\begin{quote}
\emph{śahernā kendramā̃ moṭũ bajār che.}
\end{quote}

\begin{quote}
There is a big market in the centre of the city.
\end{quote}

A neuter noun.

\begin{cousinweb}
\textbf{\gu{કેન્દ્ર}} (\emph{kendra}) is a tatsama, Sanskrit \emph{kendra}, \textquotedblleft{}the centre of a circle\textquotedblright{}, a word Sanskrit itself took from Greek (Monier-Williams).
\end{cousinweb}

\begin{grammarlens}[title={What you've built}]
One more word for getting around.
\end{grammarlens}

\subsection*{Your turn}
\begin{itemize}
\raggedright
  \item \emph{Say it:} \emph{kendra}
  \item \emph{Say it:} \emph{kendra}, once more
  \item \emph{Say it:} \emph{śahernā kendramā̃ moṭũ bajār che.} — There is a big market in the centre of the city.
  \item \emph{Recall:} say the Gujarati for a clock tower, then the Gujarati for to hunt for, then say \emph{kendra} again
\end{itemize}

\begin{tcolorbox}[breakable,colback=teal!4,colframe=teal!35!black,title={Before you move on}]
What does \gu{કેન્દ્ર} mean? (\textquotedblleft{}A centre\textquotedblright{}.) Say it once more.
\end{tcolorbox}
\section[\texorpdfstring{\gu{વિસ્તાર} (vistār)}{વિસ્તાર (vistār)}]{\gu{વિસ્તાર} (vistār) — an area; a part of a town}

\label{lesson:GU-C416-vistar}



\begin{quote}
Before the new one: say the Gujarati for to hunt for, then the Gujarati for a centre.
\end{quote}

\subsection*{You'll want to know: \gu{વિસ્તાર}}
\textbf{\gu{વિસ્તાર}} — \emph{vistār} — \textquotedblleft{}an area; a part of a town\textquotedblright{}.

\begin{quote}
\gu{આ} \gu{વિસ્તારમાં} \gu{ઘણી} \gu{દુકાનો} \gu{છે}.
\end{quote}

\begin{quote}
\emph{ā vistārmā̃ ghaṇī dukāno che.}
\end{quote}

\begin{quote}
There are many shops in this area.
\end{quote}

A masculine noun.

\begin{cousinweb}
\textbf{\gu{વિસ્તાર}} (\emph{vistār}) is a tatsama, Sanskrit \emph{vistāra}, \textquotedblleft{}spreading, extent\textquotedblright{} (Monier-Williams): the stretch of ground a place spreads over.
\end{cousinweb}

\begin{grammarlens}[title={What you've built}]
One more word for getting around.
\end{grammarlens}

\subsection*{Your turn}
\begin{itemize}
\raggedright
  \item \emph{Say it:} \emph{vistār}
  \item \emph{Say it:} \emph{vistār}, once more
  \item \emph{Say it:} \emph{ā vistārmā̃ ghaṇī dukāno che.} — There are many shops in this area.
  \item \emph{Recall:} say the Gujarati for to hunt for, then the Gujarati for a centre, then say \emph{vistār} again
\end{itemize}

\begin{tcolorbox}[breakable,colback=teal!4,colframe=teal!35!black,title={Before you move on}]
What does \gu{વિસ્તાર} mean? (\textquotedblleft{}An area; a part of a town\textquotedblright{}.) Say it once more.
\end{tcolorbox}
\section[\texorpdfstring{\gu{પરું} (parũ)}{પરું (parũ)}]{\gu{પરું} (parũ) — a suburb; an outlying quarter}

\label{lesson:GU-C416-paru}



\begin{quote}
Before the new one: say the Gujarati for a centre, then the Gujarati for an area.
\end{quote}

\subsection*{You'll want to know: \gu{પરું}}
\textbf{\gu{પરું}} — \emph{parũ} — \textquotedblleft{}a suburb; an outlying quarter\textquotedblright{}.

\begin{quote}
\gu{અમારું} \gu{ઘર} \gu{શહેરના} \gu{પરામાં} \gu{છે}.
\end{quote}

\begin{quote}
\emph{amārũ ghar śahernā parāmā̃ che.}
\end{quote}

\begin{quote}
Our house is in a suburb of the city.
\end{quote}

A neuter noun.

\begin{cousinweb}
\textbf{\gu{પરું}} (\emph{parũ}) is a tadbhava of Sanskrit \emph{pura}, \textquotedblleft{}a town, a fortress\textquotedblright{}; CDIAL lists the Gujarati word, \textquotedblleft{}a suburb\textquotedblright{} (CDIAL 8278).
\end{cousinweb}

\begin{grammarlens}[title={What you've built}]
One more word for getting around.
\end{grammarlens}

\subsection*{Your turn}
\begin{itemize}
\raggedright
  \item \emph{Say it:} \emph{parũ}
  \item \emph{Say it:} \emph{parũ}, once more
  \item \emph{Say it:} \emph{amārũ ghar śahernā parāmā̃ che.} — Our house is in a suburb of the city.
  \item \emph{Recall:} say the Gujarati for a centre, then the Gujarati for an area, then say \emph{parũ} again
\end{itemize}

\begin{tcolorbox}[breakable,colback=teal!4,colframe=teal!35!black,title={Before you move on}]
What does \gu{પરું} mean? (\textquotedblleft{}A suburb; an outlying quarter\textquotedblright{}.) Say it once more.
\end{tcolorbox}
\section[\texorpdfstring{\gu{ભાગોળ} (bhāgoḷ)}{ભાગોળ (bhāgoḷ)}]{\gu{ભાગોળ} (bhāgoḷ) — the outskirts of a village (by its gate)}

\label{lesson:GU-C416-bhagol}



\begin{quote}
Before the new one: say the Gujarati for an area, then the Gujarati for a suburb.
\end{quote}

\subsection*{You'll want to know: \gu{ભાગોળ}}
\textbf{\gu{ભાગોળ}} — \emph{bhāgoḷ} — \textquotedblleft{}the outskirts of a village (by its gate)\textquotedblright{}.

\begin{quote}
\gu{સાંજે} \gu{ગામની} \gu{ભાગોળે} \gu{બાળકો} \gu{રમે} \gu{છે}.
\end{quote}

\begin{quote}
\emph{sā̃je gāmnī bhāgoḷe bāḷko rame che.}
\end{quote}

\begin{quote}
In the evening children play at the edge of the village.
\end{quote}

A feminine noun.

\begin{cousinweb}
\textbf{\gu{ભાગોળ}} (\emph{bhāgoḷ}) goes back to \emph{*bhagnatōlī}, a form CDIAL reconstructs; CDIAL lists the Gujarati word for the gate in a village wall and the ground around it (CDIAL 9352).
\end{cousinweb}

\begin{grammarlens}[title={What you've built}]
One more word for getting around.
\end{grammarlens}

\subsection*{Your turn}
\begin{itemize}
\raggedright
  \item \emph{Say it:} \emph{bhāgoḷ}
  \item \emph{Say it:} \emph{bhāgoḷ}, once more
  \item \emph{Say it:} \emph{sā̃je gāmnī bhāgoḷe bāḷko rame che.} — In the evening children play at the edge of the village.
  \item \emph{Recall:} say the Gujarati for an area, then the Gujarati for a suburb, then say \emph{bhāgoḷ} again
\end{itemize}

\begin{tcolorbox}[breakable,colback=teal!4,colframe=teal!35!black,title={Before you move on}]
What does \gu{ભાગોળ} mean? (\textquotedblleft{}The outskirts of a village (by its gate)\textquotedblright{}.) Say it once more.
\end{tcolorbox}
\section[\texorpdfstring{\gu{ભમવું} (bhamvũ)}{ભમવું (bhamvũ)}]{\gu{ભમવું} (bhamvũ) — to roam; to travel about}

\label{lesson:GU-C416-bhamvu}



\begin{quote}
Before the new one: say the Gujarati for a suburb, then the Gujarati for the outskirts of a village.
\end{quote}

\subsection*{You'll want to know: \gu{ભમવું}}
\textbf{\gu{ભમવું}} — \emph{bhamvũ} — \textquotedblleft{}to roam; to travel about\textquotedblright{}.

\begin{quote}
\gu{તે} \gu{ઘણાં} \gu{શહેરોમાં} \gu{ભમ્યો}.
\end{quote}

\begin{quote}
\emph{te ghaṇā̃ śaheromā̃ bhamyo.}
\end{quote}

\begin{quote}
He roamed about many cities.
\end{quote}

Past (he …): \gu{તે} \gu{ઘણું} \gu{ભમ્યો}, \emph{te ghaṇũ bhamyo}.

\begin{cousinweb}
\textbf{\gu{ભમવું}} (\emph{bhamvũ}) is a tadbhava of Sanskrit \emph{bhramati}, \textquotedblleft{}wanders, roams\textquotedblright{}; CDIAL lists the Gujarati verb (CDIAL 9648).
\end{cousinweb}

\begin{grammarlens}[title={What you've built}]
One more word for getting around.
\end{grammarlens}

\subsection*{Your turn}
\begin{itemize}
\raggedright
  \item \emph{Say it:} \emph{bhamvũ}
  \item \emph{Say it:} \emph{bhamvũ}, once more
  \item \emph{Say it:} \emph{te ghaṇā̃ śaheromā̃ bhamyo.} — He roamed about many cities.
  \item \emph{Recall:} say the Gujarati for a suburb, then the Gujarati for the outskirts of a village, then say \emph{bhamvũ} again
\end{itemize}

\begin{tcolorbox}[breakable,colback=teal!4,colframe=teal!35!black,title={Before you move on}]
What does \gu{ભમવું} mean? (\textquotedblleft{}To roam; to travel about\textquotedblright{}.) Say it once more.
\end{tcolorbox}
\section[(dialogue)]{Looking back: from Lodgings to Areas and Places}

\label{lesson:GU-R416-review-409-416}



\begin{quote}
Say the Gujarati for hospitality and a manager.
\end{quote}

\subsection*{Your turn}
Say these aloud:

\begin{itemize}
\raggedright
  \item hospitality, a manager, a convenience, an inconvenience, to welcome
  \item a bathroom, a toilet, a water heater, to bail out, to wring out
  \item a quilt, a cushion for the head, a mosquito net, to roll up, to put to sleep
  \item a tent, a halt for the night, a night's stay, to drive in, to spend (time)
  \item towards, a route, the far side, to lead, to point out
  \item the east, the west, the north, the wrong way round, a compass
  \item a crossroads, a side lane, a landmark, a clock tower, to hunt for
  \item a centre, an area, a suburb, the outskirts of a village, to travel about
\end{itemize}

\begin{tcolorbox}[breakable,colback=teal!4,colframe=teal!35!black,title={Before you move on}]
Hospitality? (\textbf{\gu{પરોણાગત}}, \emph{paroṇāgat}.) A manager? (\textbf{\gu{મેનેજર}}, \emph{menejar}.) A convenience? (\textbf{\gu{સગવડ}}, \emph{sagvaḍ}.) An inconvenience? (\textbf{\gu{અગવડ}}, \emph{agvaḍ}.) To welcome? (\textbf{\gu{આવકારવું}}, \emph{āvkārvũ}.) A bathroom? (\textbf{\gu{બાથરૂમ}}, \emph{bāthrūm}.) A toilet? (\textbf{\gu{સંડાસ}}, \emph{sanḍās}.) A water heater? (\textbf{\gu{ગીઝર}}, \emph{gījhar}.) To bail out? (\textbf{\gu{ઉલેચવું}}, \emph{ulecvũ}.) To wring out? (\textbf{\gu{નિચોવવું}}, \emph{nicovvũ}.) A quilt? (\textbf{\gu{રજાઈ}}, \emph{rajāī}.) A cushion for the head? (\textbf{\gu{ઓશીકું}}, \emph{ośīkũ}.) A mosquito net? (\textbf{\gu{મચ્છરદાની}}, \emph{macchardānī}.) To roll up? (\textbf{\gu{સંકેલવું}}, \emph{sankelvũ}.) To put to sleep? (\textbf{\gu{સુવડાવવું}}, \emph{suvḍāvvũ}.) A tent? (\textbf{\gu{ડેરો}}, \emph{ḍero}.) A halt for the night? (\textbf{\gu{પડાવ}}, \emph{paḍāv}.) A night's stay? (\textbf{\gu{રાતવાસો}}, \emph{rātvāso}.) To drive in? (\textbf{\gu{ખોડવું}}, \emph{khoḍvũ}.) To spend (time)? (\textbf{\gu{ગાળવું}}, \emph{gāḷvũ}.) Towards? (\textbf{\gu{તરફ}}, \emph{taraph}.) A route? (\textbf{\gu{માર્ગ}}, \emph{mārg}.) The far side? (\textbf{\gu{પાર}}, \emph{pār}.) To lead? (\textbf{\gu{દોરવું}}, \emph{dorvũ}.) To point out? (\textbf{\gu{ચીંધવું}}, \emph{cī̃dhvũ}.) The east? (\textbf{\gu{પૂર્વ}}, \emph{pūrva}.) The west? (\textbf{\gu{પશ્ચિમ}}, \emph{paścim}.) The north? (\textbf{\gu{ઉત્તર}}, \emph{uttar}.) The wrong way round? (\textbf{\gu{ઊલટું}}, \emph{ūlṭũ}.) A compass? (\textbf{\gu{હોકાયંત્ર}}, \emph{hokāyantra}.) A crossroads? (\textbf{\gu{ચોકડી}}, \emph{cokḍī}.) A side lane? (\textbf{\gu{ખાંચો}}, \emph{khā̃co}.) A landmark? (\textbf{\gu{નિશાની}}, \emph{niśānī}.) A clock tower? (\textbf{\gu{ટાવર}}, \emph{ṭāvar}.) To hunt for? (\textbf{\gu{ગોતવું}}, \emph{gotvũ}.) A centre? (\textbf{\gu{કેન્દ્ર}}, \emph{kendra}.) An area? (\textbf{\gu{વિસ્તાર}}, \emph{vistār}.) A suburb? (\textbf{\gu{પરું}}, \emph{parũ}.) The outskirts of a village? (\textbf{\gu{ભાગોળ}}, \emph{bhāgoḷ}.) To travel about? (\textbf{\gu{ભમવું}}, \emph{bhamvũ}.)
\end{tcolorbox}
